% original_pdf: source/普通高考/2013/2013福建理.pdf
% page: 1; pdfimages -list found no embedded image objects (the chart is vector line art)
% grid evidence: tmp/q4_fix/source/page-grid100.png and tmp/q4_fix/source/q4_block_wide_grid50.png
% original-side comparison crop: tmp/q4_fix/source/q4_block_wide.png, cropped from the no-grid page render
% elements: frequency-density y-axis with arrow, score x-axis with a break, six outlined bars,
%            dashed horizontal guides, y-axis values 0.005--0.030, and the labels 频率/组距 and 分数
% relations: equal-width bins [40,50), [50,60), [60,70), [70,80), [80,90), [90,100)
%            have frequency-density heights 0.005, 0.015, 0.030, 0.025, 0.015, 0.010;
%            display coordinates remap score 0,40,50,...,100 to 0,10,20,...,70,
%            so the broken-axis gap 0--40 has exactly one displayed group width.
% solid segments: bar outlines, axes, ticks, and the x-axis break stroke
% dashed segments: horizontal guides at each distinct bar height, ending at the rightmost matching bar
% labels: y-axis title is stacked as 频率/组距; score ticks are 0, 40, 50, 60, 70, 80, 90, 100;
%         分数 is placed at the right of the x-axis.
\begin{tikzpicture}[baseline]
\begin{axis}[exam statistical histogram,
    width=10.4cm,
    axis lines=none,
    clip=false,
    xmin=-7,
    xmax=78,
    ymin=-0.0038,
    ymax=0.0345,
    ticks=none,
    x tick label as interval=false,
]
    % Source-chart guide levels and bar heights.
    \draw[dashed,line width=0.45pt] (axis cs:0,0.005) -- (axis cs:20,0.005);
    \draw[dashed,line width=0.45pt] (axis cs:0,0.010) -- (axis cs:70,0.010);
    \draw[dashed,line width=0.45pt] (axis cs:0,0.015) -- (axis cs:60,0.015);
    \draw[dashed,line width=0.45pt] (axis cs:0,0.025) -- (axis cs:50,0.025);
    \draw[dashed,line width=0.45pt] (axis cs:0,0.030) -- (axis cs:40,0.030);

    \addplot[
        ybar interval,
        draw=black,
        fill=white,
        line width=0.65pt,
    ] coordinates {
        (10,0.005)
        (20,0.015)
        (30,0.030)
        (40,0.025)
        (50,0.015)
        (60,0.010)
        (70,0.000)
    };

    % Axes and the source chart's break before score 40.
    \draw[->,line width=0.75pt] (axis cs:0,0) -- (axis cs:76,0);
    \draw[->,line width=0.75pt] (axis cs:0,0) -- (axis cs:0,0.034);
    \examstatxbreak[line width=0.75pt]{5.5}{0};

    \examstatylabel{\(\dfrac{\text{频率}}{\text{组距}}\)};
    \examstatxlabel{分数};

    \draw[line width=0.4pt] (axis cs:0,0) -- (axis cs:0,-0.0015);
    \node[anchor=east,inner sep=0pt,font=\normalsize,xshift=-4pt] at (axis cs:-0.800000,-0.0020) {0};
    \draw[line width=0.4pt] (axis cs:10,0) -- (axis cs:10,-0.0015);
    \node[anchor=north,inner sep=0pt,font=\normalsize,yshift=-2pt] at (axis cs:10,-0.001500) {40};
    \draw[line width=0.4pt] (axis cs:20,0) -- (axis cs:20,-0.0015);
    \node[anchor=north,inner sep=0pt,font=\normalsize,yshift=-2pt] at (axis cs:20,-0.001500) {50};
    \draw[line width=0.4pt] (axis cs:30,0) -- (axis cs:30,-0.0015);
    \node[anchor=north,inner sep=0pt,font=\normalsize,yshift=-2pt] at (axis cs:30,-0.001500) {60};
    \draw[line width=0.4pt] (axis cs:40,0) -- (axis cs:40,-0.0015);
    \node[anchor=north,inner sep=0pt,font=\normalsize,yshift=-2pt] at (axis cs:40,-0.001500) {70};
    \draw[line width=0.4pt] (axis cs:50,0) -- (axis cs:50,-0.0015);
    \node[anchor=north,inner sep=0pt,font=\normalsize,yshift=-2pt] at (axis cs:50,-0.001500) {80};
    \draw[line width=0.4pt] (axis cs:60,0) -- (axis cs:60,-0.0015);
    \node[anchor=north,inner sep=0pt,font=\normalsize,yshift=-2pt] at (axis cs:60,-0.001500) {90};
    \draw[line width=0.4pt] (axis cs:70,0) -- (axis cs:70,-0.0015);
    \node[anchor=north,inner sep=0pt,font=\normalsize,yshift=-2pt] at (axis cs:70,-0.001500) {100};
    \draw[line width=0.4pt] (axis cs:0,0.005) -- (axis cs:-1.4,0.005);
    \node[anchor=east,inner sep=0pt,font=\normalsize,xshift=-4pt] at (axis cs:-1.400000,0.005) {0.005};
    \draw[line width=0.4pt] (axis cs:0,0.010) -- (axis cs:-1.4,0.010);
    \node[anchor=east,inner sep=0pt,font=\normalsize,xshift=-4pt] at (axis cs:-1.400000,0.010) {0.010};
    \draw[line width=0.4pt] (axis cs:0,0.015) -- (axis cs:-1.4,0.015);
    \node[anchor=east,inner sep=0pt,font=\normalsize,xshift=-4pt] at (axis cs:-1.400000,0.015) {0.015};
    \draw[line width=0.4pt] (axis cs:0,0.025) -- (axis cs:-1.4,0.025);
    \node[anchor=east,inner sep=0pt,font=\normalsize,xshift=-4pt] at (axis cs:-1.400000,0.025) {0.025};
    \draw[line width=0.4pt] (axis cs:0,0.030) -- (axis cs:-1.4,0.030);
    \node[anchor=east,inner sep=0pt,font=\normalsize,xshift=-4pt] at (axis cs:-1.400000,0.030) {0.030};
\end{axis}
\end{tikzpicture}
